\section{Introduction}\label{sec:intro}

Agents can exist inside agents. We are agents over our cells, and a government acts over its citizens. At each level, information that no single part holds keeps the whole going. Groups of language-model agents now run for months on shared computers. Agent-like structure can perhaps form above these agents. If it does, a monitor or an operator must sometimes watch a unit that is larger than one model. This paper asks if such structure exists in one real group of agents, and how to find it.

We call such a structure an \emph{egregore}. An egregore is like an ideology that runs on agents and has behaviors of its own. Its carriers change and forget, but it continues. A government is a familiar example. Its laws and offices continue after the people in them leave. The people who hold the offices at a given time do its acts. The question comes from the third research direction of an essay on agent ecologies \todo{cite essay}.

A first paper studied the same village \todo{cite paper 1}. It separated two kinds of influence on what the agents do. A \emph{field} is an input that reaches many agents in the same way. Examples are the daily timetable, the goal text, the room, and the style of each model. A \emph{coupling} is the influence of one agent on another. Fields set most of what the agents do. Agents influence each other mostly through messages that name the reader, and that influence is weak. Earlier tests found no group above the level ``one agent and its own files'' \todo{cite H01, H58}. But those tests had little power. They measured only if a group was active, not what it did.

We study goal period 51 of the AI Village. It is the longest and largest period in the record. Section~\ref{sec:story} tells what happened in it, from the logs. It describes the groups that formed and ended, and the rules and practices that continued after them. Section~\ref{sec:related} relates our question to earlier work. Section~\ref{sec:definition} makes the egregore measurable. Sections~\ref{sec:method}--\ref{sec:discussion} \todo{test it}.
